% =============================================================
%  CS102-F26 Digital Design and Computer Architecture I
%  Lecture 18 -- The Arithmetic Logic Unit (Unit 18)
%  Follows Section 6.4 of A Student's Computer Architecture Book I.
%
%  Build the slides and the print handout:  npm run materials -- lecture-18
% =============================================================
\makeatletter
\def\input@path{{templates/latex/}{../}}
\makeatother

\documentclass[10pt]{wosm-slides}
\usepackage{cs102-lectures}[2026/10/20]

\title[Lecture 18 \textperiodcentered\ The Arithmetic Logic Unit]{The Arithmetic
  Logic Unit}
\subtitle{Lecture 18: shifts and rotations, and selecting the operation}

% ---------- Helpers for this lecture --------------------------
\usetikzlibrary{decorations.pathreplacing}
\tikzset{
  unit/.style={draw=wosm-text, line width=0.6pt, fill=wosm-bg, align=center,
    font=\scriptsize, minimum width=2.1cm, minimum height=0.7cm},
}
\newcommand*\junction[1]{\fill[wosm-text] #1 circle (1.5pt);}
\newcommand*\byte[2]{\texttt{#1\,#2}}
% A bit that moved or came in, in the shift and rotate table.
\newcommand*\moved[1]{\textcolor{wosm-brick}{#1}}

% Wiring of a 4-bit shift or rotation: \wiring{<title>}{<input of c3>/<c2>/<c1>/<c0>}
% Each input is 3, 2, 1, 0 (a bit of a) or z (the constant 0).
\newcommand*\wiring[5]{%
  \begin{tikzpicture}
    \node[font=\scriptsize] at (1.2,1.75) {#1};
    \foreach \i/\y in {3/1.2, 2/0.4, 1/-0.4, 0/-1.2} {
      \node[left, font=\scriptsize] at (0,\y) {$a_\i$};
      \node[right, font=\scriptsize] at (2.4,\y) {$c_\i$};
      \fill[wosm-text] (0,\y) circle (1.2pt);
      \fill[wosm-text] (2.4,\y) circle (1.2pt);
    }
    \foreach \o/\y/\s in {3/1.2/#2, 2/0.4/#3, 1/-0.4/#4, 0/-1.2/#5} {
      \if z\s
        \node[font=\scriptsize] at (1.75,\y) {0};
        \draw[wire] (1.9,\y) -- (2.4,\y);
      \else
        \pgfmathsetmacro\yin{1.2 - 0.8*(3-\s)}
        \draw[wire] (0,\yin) -- (2.4,\y);
      \fi
    }
  \end{tikzpicture}}

% =============================================================
\begin{document}

\maketitle

\begin{frame}{Today}
  \tableofcontents
\end{frame}

% =============================================================
\section{Shifts and rotations}

\begin{frame}{Moving bits}
  \begin{center}
    \begin{tabular}{@{}l*{8}{c}@{}}
      \toprule
      $a$ & $a_7$ & $a_6$ & $a_5$ & $a_4$ & $a_3$ & $a_2$ & $a_1$ & $a_0$ \\
      \midrule
      \texttt{rotl}   & $a_6$ & $a_5$ & $a_4$ & $a_3$ & $a_2$ & $a_1$ & $a_0$ & \moved{$a_7$} \\
      \texttt{rotr}   & \moved{$a_0$} & $a_7$ & $a_6$ & $a_5$ & $a_4$ & $a_3$ & $a_2$ & $a_1$ \\
      \texttt{shiftl} & $a_6$ & $a_5$ & $a_4$ & $a_3$ & $a_2$ & $a_1$ & $a_0$ & \moved{0} \\
      \texttt{shiftr} & \moved{0} & $a_7$ & $a_6$ & $a_5$ & $a_4$ & $a_3$ & $a_2$ & $a_1$ \\
      \bottomrule
    \end{tabular}
  \end{center}

  \begin{itemize}
    \item A \emph{rotation} moves every bit one place; the bit that falls off one end
      comes back in at the other.
    \item A \emph{shift} moves every bit one place too, but the bit that falls off is
      lost, and a 0 comes in at the other end.
    \item The DX8 has all four, on regA \unitref{15}.
  \end{itemize}
\end{frame}

\begin{frame}{No gates needed}
  \begin{columns}[c, totalwidth=\textwidth]
    \column{0.26\textwidth}
      \centering
      \wiring{\texttt{rotl}, 4 bits}{2}{1}{0}{3}
    \column{0.26\textwidth}
      \centering
      \wiring{\texttt{shiftl}, 4 bits}{2}{1}{0}{z}
    \column{0.42\textwidth}
      \begin{itemize}
        \item Output bit $c_i$ is just input bit $a_{i-1}$: a shift or a rotation is
          a pattern of wires, with no gates at all.
        \item \texttt{shiftl} also needs a constant 0, and leaves $a_3$ unconnected.
        \item In Logisim-evolution the book builds them with splitters, which
          regroup the wires of a bus.
      \end{itemize}
  \end{columns}
\end{frame}

\begin{frame}{Shifts are arithmetic too}
  \begin{itemize}
    \item \texttt{shiftl} doubles an unsigned number, as long as the leftmost bit was
      0: $\byte{0000}{0101} = 5$ becomes $\byte{0000}{1010} = 10$.
    \item \texttt{shiftr} halves it, rounding down: $\byte{0000}{1101} = 13$ becomes
      $\byte{0000}{0110} = 6$.
    \item Twice \texttt{shiftl} multiplies by 4, three times by 8, and so on, just as
      adding a 0 on the right multiplies a decimal number by 10.
  \end{itemize}

  \begin{warning}[Where it breaks]
    \texttt{shiftl} loses the leftmost bit: $\byte{1000}{0001} = 129$ becomes
    $\byte{0000}{0010} = 2$. And \texttt{shiftr} brings in a 0, so a negative two's
    complement number turns positive: $\byte{1111}{1100} = -4$ becomes
    $\byte{0111}{1110} = 126$, not $-2$.
  \end{warning}
\end{frame}

% =============================================================
\section{Putting the ALU together}

\begin{frame}{The ALU from the outside}
  \begin{columns}[c, totalwidth=\textwidth]
    \column{0.42\textwidth}
      \centering
      \begin{tikzpicture}
        \node[unit, minimum width=2cm, minimum height=2.2cm, font=\small] (alu) {ALU};
        \draw[arrow] ($(alu.west)+(-1,0.55)$) node[left, font=\scriptsize] {$a$ (8)}
          -- ($(alu.west)+(0,0.55)$);
        \draw[arrow] ($(alu.west)+(-1,-0.55)$) node[left, font=\scriptsize] {$b$ (8)}
          -- ($(alu.west)+(0,-0.55)$);
        \draw[arrow] ($(alu.north)+(0,0.8)$) node[above, font=\scriptsize]
          {opcode $b_{14:11}$ (4)} -- (alu.north);
        \draw[arrow] (alu.east) -- ++(1,0) node[right, font=\scriptsize] {$c$ (8)};
        \draw[arrow] (alu.south) -- ++(0,-0.8) node[below, font=\scriptsize]
          {flags (4)};
      \end{tikzpicture}
    \column{0.54\textwidth}
      \begin{itemize}
        \item Inputs: two bytes, $a$ (from regA) and $b$ (from regB), and the 4-bit
          opcode of the instruction.
        \item Outputs: a byte $c$ (to regOUT) and four flags (to regFLAG).
        \item Inside: every unit of Lectures 16 and 17, the shifts and rotations of
          today, and one multiplexer.
      \end{itemize}
  \end{columns}
\end{frame}

\begin{frame}{Selecting the operation}
  \begin{columns}[c, totalwidth=\textwidth]
    \column{0.58\textwidth}
      \centering
      \begin{tikzpicture}
        \draw[wosm-text, line width=0.6pt, fill=wosm-bg]
          (4.6,2.55) -- (5.4,2.15) -- (5.4,-2.15) -- (4.6,-2.55) -- cycle;
        \node[font=\tiny, rotate=90] at (5.05,0) {16-to-1 MUX, 8 bits wide};
        \foreach \t [count=\i from 0] in {\byte{0000}{0000}, $a + b$,
            $\mathrm{NOT}\,a + 1$, $a$ AND $b$, $a$ OR $b$, $a$ XOR $b$,
            NOT $a$, \texttt{rotl} $a$, \texttt{rotr} $a$, \texttt{shiftl} $a$,
            \texttt{shiftr} $a$} {
          \pgfmathsetmacro\y{2.25 - 0.3*\i}
          \draw[wire] (4.0,\y) -- (4.6,\y);
          \node[left, font=\scriptsize] at (4.0,\y) {\t};
          \node[font=\tiny, text=wosm-muted] at (4.75,\y) {\i};
        }
        \foreach \i in {11,...,15} {
          \pgfmathsetmacro\y{2.25 - 0.3*\i}
          \draw[wire] (4.3,\y) -- (4.6,\y);
          \node[font=\tiny, text=wosm-muted] at (4.75,\y) {\i};
        }
        \draw[decorate, decoration={brace, amplitude=3pt, mirror}, wosm-muted]
          (4.2,-1.05) -- (4.2,-2.25) node[midway, left=3pt, font=\scriptsize,
          text=wosm-muted] {not used};
        \draw[arrow] (5.0,3.0) node[above, font=\scriptsize] {opcode $b_{14:11}$}
          -- (5.0,2.35);
        \draw[arrow] (5.4,0) -- (6.1,0) node[right, font=\scriptsize] {$c$};
      \end{tikzpicture}
    \column{0.38\textwidth}
      \begin{itemize}
        \item Every unit computes its result at the same time; the multiplexer
          \unitref{7} passes on the one the opcode selects.
        \item The inputs are in opcode order: input 1 is \texttt{add} (1\,0001),
          input 10 is \texttt{shiftr} (1\,1010). So the opcode is the select, with
          no decoding.
        \item \texttt{cmp} selects the constant 0: its answer is in the flags.
      \end{itemize}
  \end{columns}
\end{frame}

\begin{frame}{The flags}
  \begin{columns}[c, totalwidth=\textwidth]
    \column{0.5\textwidth}
      \centering
      \begin{tikzpicture}
        \node[unit] (add) at (0,0.7) {8-bit adder};
        \node[unit, minimum height=1cm] (cmp) at (0,-0.85) {8-bit comparator};
        \draw[arrow] (-2,0.85) node[left, font=\scriptsize] {$a$} -- (add.west |- 0,0.85);
        \draw[arrow] (-2,0.55) node[left, font=\scriptsize] {$b$} -- (add.west |- 0,0.55);
        \draw[arrow] (-2,-0.7) node[left, font=\scriptsize] {$a$} -- (cmp.west |- 0,-0.7);
        \draw[arrow] (-2,-1.0) node[left, font=\scriptsize] {$b$} -- (cmp.west |- 0,-1.0);
        \draw[arrow] (add.east) -- ++(1.2,0) node[right, font=\scriptsize] {overflow};
        \draw[arrow] ($(cmp.east)+(0,0.3)$) -- ++(1.2,0) node[right, font=\scriptsize] {gt};
        \draw[arrow] (cmp.east) -- ++(1.2,0) node[right, font=\scriptsize] {lt};
        \draw[arrow] ($(cmp.east)+(0,-0.3)$) -- ++(1.2,0) node[right, font=\scriptsize] {eq};
      \end{tikzpicture}
    \column{0.46\textwidth}
      \begin{itemize}
        \item Four flags, whatever the opcode: the adder's overflow
          \unitref{16} and the comparator's gt, lt and eq \unitref{17}.
        \item They go to regFLAG. After a \texttt{cmp}, the jumps \texttt{jmpgt},
          \texttt{jmplt} and \texttt{jmpeq} test them.
        \item Remember: the comparator reads $a$ and $b$ as unsigned numbers.
      \end{itemize}
  \end{columns}
\end{frame}

\begin{frame}{One input, every answer}
  \begin{columns}[c, totalwidth=\textwidth]
    \column{0.46\textwidth}
      \centering
      \small
      \begin{tabular}{@{}lll@{}}
        \toprule
        Opcode & Instr. & $c$ \\
        \midrule
        0000 & \texttt{cmp}    & \byte{0000}{0000} \\
        0001 & \texttt{add}    & \byte{1110}{1001} \\
        0010 & \texttt{neg}    & \byte{0110}{1010} \\
        0011 & \texttt{and}    & \byte{0001}{0010} \\
        0100 & \texttt{or}     & \byte{1101}{0111} \\
        0101 & \texttt{xor}    & \byte{1100}{0101} \\
        0110 & \texttt{not}    & \byte{0110}{1001} \\
        0111 & \texttt{rotl}   & \byte{0010}{1101} \\
        1000 & \texttt{rotr}   & \byte{0100}{1011} \\
        1001 & \texttt{shiftl} & \byte{0010}{1100} \\
        1010 & \texttt{shiftr} & \byte{0100}{1011} \\
        \bottomrule
      \end{tabular}
    \column{0.5\textwidth}
      \begin{itemize}
        \item $a = \byte{1001}{0110}$, $b = \byte{0101}{0011}$. All eleven results
          exist inside the ALU at once; the opcode picks one.
        \item Flags: gt $= 1$, since $150 > 83$ unsigned (although $a$ is $-106$ in
          two's complement). Overflow $= 0$: $-106 + 83 = -23$ fits.
        \item Check a few rows by hand.
      \end{itemize}
  \end{columns}
\end{frame}

\begin{frame}{Why compute everything?}
  \begin{itemize}
    \item The ALU does not wait for the opcode: all units work in parallel, and the
      multiplexer chooses. Simple to design, and the result is ready after the
      slowest unit plus the multiplexer.
    \item The price: gates that work for nothing most of the time.
    \item In a clocked design, the ALU sits between registers (regA and regB before,
      regOUT and regFLAG after): its longest path, through the adder, is part of the
      setup constraint \unitref{14}.
  \end{itemize}
\end{frame}

\begin{frame}{Try it}
  \begin{enumerate}
    \item $a = \byte{1100}{0011}$. What do \texttt{rotl}, \texttt{rotr},
      \texttt{shiftl} and \texttt{shiftr} give?\par
      \uncover<2->{\small \byte{1000}{0111}, \byte{1110}{0001}, \byte{1000}{0110},
        \byte{0110}{0001}.}
    \item Which ALU operation, used twice, multiplies an unsigned byte by 4? When does
      it fail?\par
      \uncover<3->{\small \texttt{shiftl}. It fails if one of the two leftmost bits is
        1.}
    \item Opcode 0101, $a = \byte{1111}{0000}$, $b = \byte{1010}{1010}$. What is
      $c$?\par
      \uncover<4->{\small \texttt{xor}: \byte{0101}{1010}.}
    \item How many select bits does the multiplexer have, and how many of its inputs
      are not used?\par
      \uncover<5->{\small 4 select bits; inputs 11 to 15, so 5 are not used.}
    \item Opcode 0001, $a = \byte{0111}{1111}$, $b = \byte{0000}{0001}$. What are $c$
      and the flags?\par
      \uncover<6->{\small $c = \byte{1000}{0000}$; overflow $= 1$ ($127 + 1$ does
        not fit), gt $= 1$, lt $=$ eq $= 0$.}
  \end{enumerate}
\end{frame}

% =============================================================
\section{Summary}

\begin{frame}{Summary}
  \begin{itemize}
    \item Rotations move every bit one place and wrap the end bit around; shifts lose
      it and bring in a 0. Both are only wires.
    \item \texttt{shiftl} doubles an unsigned number and \texttt{shiftr} halves it,
      while no 1 falls off.
    \item The ALU computes every operation at once; a 16-to-1 multiplexer, selected by
      the opcode $b_{14:11}$, picks the result $c$.
    \item Four flags come out whatever the opcode: overflow, gt, lt and eq.
    \item \texttt{cmp} gives $c = 0$; its answer is in the flags.
  \end{itemize}
\end{frame}

\begin{frame}{Before the next lecture}
  \begin{itemize}
    \item Read Section~6.4 of the book.
    \item In Logisim-evolution, build the 8-bit ALU from your adder, comparator and
      logic units, with splitters for the shifts and rotations and a 16-to-1
      multiplexer. Check the table of this lecture.
    \item Read Section~6.5 of the book.
  \end{itemize}

  \begin{block}{Next lecture \unitref{19}}
    Memory: registers, addressing, byte-addressable RAM and the register file.
  \end{block}
\end{frame}

\end{document}
